\begin{helper}[Decreasing restoration trace for Algorithm (A)]
Let \(L^*=\noderef{ExtensionBaseRestorationIndices}\). Suppose that every
slot \(t\in L^*\) has exact certified values
\[
  v_t,\gamma_t,\delta_t,\qquad
  v(t+1)=v_t,\quad \gamma(t+1)=\gamma_t,\quad
  \delta(t+1)=\delta_t.
\]
For any starting coloring \(\sigma_n\), there exists an
\noderef{ExtensionBaseRestorationTrace} based at \(\sigma_n\).
\end{helper}
\begin{proof}
The set \(L^*\) is a finite set of finite ordinals. Let \(r=|L^*|\), and
list its members in the unique strictly decreasing order
\[
                         t_0>t_1>\cdots>t_{r-1}.
\]
This includes \(r=0\), when the list has no entries. For each \(a<r\),
apply the displayed hypothesis to \(t_a\), and choose one triple
\((v_{t_a},\gamma_{t_a},\delta_{t_a})\) satisfying all three exact field
equalities. Define colorings recursively by
\[
 \sigma^0=\sigma_n,\qquad
 \sigma^{a+1}=\noderef{TwoColorComponentSwap}
   (\sigma^a,\gamma_{t_a},\delta_{t_a},v_{t_a}).
\]
Take the number of steps to be \(r\), the slot at position \(a\) to be
\(t_a\), the support and two colors to be the chosen triple, and the coloring
at position \(a\) to be \(\sigma^a\). The initial-coloring and transition
clauses of \noderef{ExtensionBaseRestorationTrace} hold by these recursive
definitions. The image of the slot map is exactly \(L^*\), because the
ordered list contains every member exactly once. Its strict-decrease clause
holds because the list was sorted in decreasing order. Finally, the three
certified-field clauses are precisely the properties of the chosen triple.
Thus all fields of the required trace are defined. When \(r=0\), the same
construction consists only of \(\sigma^0=\sigma_n\), and all slot-indexed
clauses are vacuous.
\end{proof}
